% ===========================================================================
%  附录 A：INDI
% ===========================================================================

\section{增量非线性动态逆的适配条件}
\label{sec:INDI}

增量非线性动态逆（INDI）利用当前角加速度和局部控制效能形成执行器增量，
可降低对完整气动模型的依赖，但其鲁棒性取决于角加速度估计、执行器同步和效能矩阵误差
\cite{sieberling2010robust,smeur2016incremental,ludena2020flight,pollack2023robust}。

设转动动力学写成
\begin{equation}
\dot{\bm\omega}=\bm f(\bm x)+\bm G(\bm x)\bm u,
\qquad
\bm G=\bm I^{-1}\bm B_M,
\label{eq:indi_plant}
\end{equation}
其中 $\bm B_M=\partial\bm M/\partial\bm u$ 的单位为
$\mathrm{N\,m}$ 对执行器单位，$\bm G$ 的单位为角加速度对执行器单位。
在当前工作点一阶展开并忽略一拍内的慢变未建模项，可得
\begin{equation}
\Delta\bm u
=\bm G_k^{-1}(\bm\nu_k-\dot{\bm\omega}_k)
=\bm B_{M,k}^{-1}\bm I
(\bm\nu_k-\dot{\bm\omega}_k).
\label{eq:indi_law}
\end{equation}
式\eqref{eq:indi_law}中的惯量只能出现一次。若直接使用力矩效能矩阵
$\bm B_M$，必须先把角加速度误差乘以 $\bm I$；若使用角加速度效能矩阵
$\bm G$，则不能再次乘或除以惯量。早期稿件把两种定义混合，产生了量纲错误。

\subsection{角加速度估计}
\label{sec:accel_estimation}

最简单的离散估计为
\begin{equation}
\widehat{\dot{\bm\omega}}_k
=\frac{\bm\omega_k-\bm\omega_{k-1}}{\Delta t},
\label{eq:bwd_euler}
\end{equation}
但微分会放大陀螺噪声。低通、固定滞后平滑、滑模微分或模型--测量互补估计都在噪声和
群延迟之间权衡\cite{levant1998robust,jiang2025adaptive,zhu2026mems}。
INDI 要求用于 $\dot{\bm\omega}_k$ 的测量和用于 $\Delta\bm u_k$ 的执行器状态在时间上对齐；
如果电机以 $\tau_m$ 滞后而控制器只使用指令值，增量模型会把执行器相位差解释为外扰。

\subsection{本平台的效能矩阵}

本平台可令 $\bm B_M=\bm B_{\mathrm{true}}$，其表达式见\cref{eq:B_true}。
由于差速列和摆角列当前使用的转速工作点并非完全同源，直接把固件矩阵移入 INDI
仍需先统一执行器状态。还需加入 $\pm15^\circ$ 摆角、$\pm0.7$ 差速、单侧转速和
增量率限制，否则累加式控制会在饱和后产生不可实现的内部状态。

\subsection{与当前固件的区别}

当前固件每拍计算
$\bm u_k=\bm B_k^{-1}\bm M_{\mathrm{cmd},k}$，不测量角加速度，也不累加
$\Delta\bm u_k$，因此不属于式\eqref{eq:indi_law}定义的 INDI。
现有 Python 对比脚本虽然采用累加式指令，却使用与固件不同的四元数误差、工作点、限幅和
执行器动态。它可作为算法原型，不能作为当前固件性能证据。

\subsection{后续验证要求}

INDI 的正式比较应与固件级联控制器共享同一非线性植物、初值、参数哈希、执行器动态和约束，
并报告角加速度估计延迟、噪声谱、饱和率和随机种子。只有满足这些条件，才可以把
跟踪误差差异归因于控制算法，而不是实现和测试条件差异。
